% $Id: responses.tex 
% !TEX root = main.tex

\newpage

\appendix

\clearpage

\subsection{Response Letter}


As required by the reviewing process, we are submitting this response letter along with a revised 
version of our paper entitled \textsl{Flik: A Back-in-time Debugger for Reinforcement Learning Programs}. For convenience, the letter is formatted as an appendix of the paper.

We would like to thank our anonymous reviewers for their efforts in reviewing our submission and 
for their useful and detailed suggestions to improve it. We have considered all remarks carefully, and 
addressed them to the best of our possibilities.

There are \thetotalremarks~ different remarks for this iteration of the paper, for which we have 
effectively answered \thetotalresponses, 
while \thetotalunsolved ~are blocked or require further feedback from the reviewers. 
Each remark 
is followed by our response and an explicit status indication:

\begin{itemize}[nosep,label=--] 
	\item \revtotal{solved} \statussolved 
	\item \revtotal{future} \statusfuture
	\item \revtotal{clarify} \statusclarify\footnote{We clarify the question made by the reviewer} 
	\item \revtotal{reconsider} \statusreconsider\footnote{We ask the reviewer to reconsider the remark in the light of arguments and clarifications we provide in the response.} 
	\item \revtotal{blocked} \statusblocked\footnote{We could not accommodate the remark in the text due to blocking constraints.} 
	\item \revtotal{feedback} \statusfeedback\footnote{We answered the remark to our best interpretation. We are unsure about what the reviewer meant.} 	
\end{itemize}

The remarks with specific changes in the paper have a reference to the location in the paper 
where the solution is found. These solutions are marked with the 
\includegraphics[width=0.03\textwidth]{./figures/fix.pdf} icon in the paper margin for easy 
identification. Additionally, in the margin we point to the specific remark(s) the modification 
is addressing.


%%
\subsection{Review \#1}



\begin{remark} \label{rem:intro}
The introduction is overly technical.
\end{remark}
\begin{solved}
We rewrote the whole introduction to have a more abstract presentation of the motivation, and reduce its technicality (\fref[vario]{par:intro}). All the technical details have been moved to a new sections providing the background on \ac{RL} (\fref{sec:rl-background} - \fref[vario]{par:background})  
\end{solved}



\endinput
